% Source: https://artofproblemsolving.com/community/c3h2360434_trapezoid_problems
% Stable identity: aops-topic-2360434-q02
Let one of the triangles have a side of length $c$. Place it as $T_1$ with vertices $(0,0)$, $(c,0)$, $(p,h)$, where $h>0$.
\begin{itemize}
\item $T_2$ has vertices $(c,0)$, $(p+c,h)$, $(p,h)$. It is the image of $T_1$ under the half-turn about the midpoint of the segment from $(c,0)$ to $(p,h)$, so it is congruent to $T_1$.
\item $T_3$ has vertices $(c,0)$, $(2c,0)$, $(p+c,h)$; it is $T_1$ moved $c$ to the right.
\end{itemize}

\begin{asy}
usepackage("amsmath");
size(6.5cm);
defaultpen(fontsize(10pt));
real c=4, p=1.3, h=2.2;
pair P0=(0,0), P1=(c,0), P2=(2c,0), Q0=(p,h), Q1=(p+c,h);
fill(P0--P1--Q0--cycle,rgb(0.85,0.92,1)); fill(P1--Q1--Q0--cycle,rgb(1,0.9,0.8)); fill(P1--P2--Q1--cycle,rgb(0.85,0.92,1));
draw(P0--P2--Q1--Q0--cycle,linewidth(0.9)); draw(P1--Q0); draw(P1--Q1);
label("$T_1$",(P0+P1+Q0)/3); label("$T_2$",(P1+Q1+Q0)/3); label("$T_3$",(P1+P2+Q1)/3);
label("$c$",(P0+P1)/2,S); label("$c$",(P1+P2)/2,S); label("$c$",(Q0+Q1)/2,N);
\end{asy}

The segments from $(c,0)$ to $(p,h)$ and to $(p+c,h)$ cut the quadrilateral with vertices $(0,0)$, $(2c,0)$, $(p+c,h)$, $(p,h)$ into exactly these three triangles, so they do not overlap and together form this quadrilateral. Its sides on $y=0$ and $y=h$ are parallel, with lengths $2c$ and $c$. The other two sides are not parallel: they point in the directions $(p,h)$ and $(p-c,h)$, which are parallel only if $c=0$. So the three triangles form a trapezoid with parallel sides $2c$ and $c$. $\blacksquare$
